一段访问控制规则里有 5 个条件，要确认两种写法是否等价，真值表得查 \(2^5=32\) 行；条件加到 10 个就是 1024 行，加到 20 个已经没人查得动。另一头，一个用来校验括号配平的正则表达式，无论怎么改都会被更深一层的嵌套打穿。两件事看上去不相干，难处却同源：能动用的表示是有限的，要覆盖的情形不是。

本章两节各处理一种``有限''。布尔代数把逻辑变成一套可以机械演算的代数，它与第一部分的命题逻辑同源，问的问题却换了——不是这句话对不对，而是算出这句话要花几个门、哪几个条件其实白写。有限自动机把无限长的历史压成有限个状态，于是可以问一个更硬的问题：有限状态能做什么，不能做什么。第二节的判断值得优先带走，它在流式处理、协议设计和输入校验里会反复被用到。

\begin{center}
\begin{tikzpicture}[x=1cm,y=1cm,font=\footnotesize,>=stealth]
  \draw[fill=blue!8] (0,0.6) rectangle (2.4,1.8);
  \node at (1.2,1.2) {组合逻辑 \(F\)};
  \draw[->] (-0.9,1.2) -- (0,1.2);
  \node[above] at (-0.45,1.25) {输入};
  \draw[->] (2.4,1.2) -- (3.4,1.2);
  \node[above] at (2.9,1.25) {输出};
  \node at (1.2,-1.3) {没有回路};
  \draw[fill=blue!8] (6,0.6) rectangle (8.4,1.8);
  \node at (7.2,1.2) {转移 \(\delta\)};
  \draw[->] (5.1,1.5) -- (6,1.5);
  \node[above] at (5.55,1.55) {输入};
  \draw[->] (8.4,1.5) -- (9.5,1.5);
  \node[above] at (8.95,1.55) {输出};
  \fill (8.9,1.5) circle (1.4pt);
  \draw[->] (8.9,1.5) -- (8.9,-0.3) -- (5.5,-0.3) -- (5.5,0.9) -- (6,0.9);
  \node[below] at (7.2,-0.35) {状态 \(q\)};
  \node at (7.2,-1.3) {有一条回线};
\end{tikzpicture}

\emph{左边没有回路，输出只由当前输入决定，这是第一节的对象；右边多出一条把状态送回输入的线，机器因此带着记忆，这是第二节的对象。整章的问题都可以挂在这两张图上：左边问一次判断要花多少门，右边问那条回线上能塞下多少信息。}
\end{center}

顺读即可，两节之间只有一处依赖：第二节用到``状态即等价类''的想法，与第一节的化简互不牵扯，从哪一节进入都行。若对命题、连接词与真值表还不熟悉，先回看\hyperref[ch:004]{``逻辑：把含混判断变成可检验命题''}会顺一些。

\hyperref[sec:04-Discrete-Mathematics/07-Boolean-Algebra-and-Automata/01]{``布尔代数与逻辑化简''}处理同一时刻的条件组合：两个表达式是否始终等价，怎样把门数降下来，以及为什么全局最优的化简是难的。它顺带回答一个工程问题——某个字段是不是根本不影响判定。\hyperref[sec:04-Discrete-Mathematics/07-Boolean-Algebra-and-Automata/02]{``有限自动机识别正则语言''}处理沿序列到达的输入：状态是对已读前缀的摘要，能合并的前缀必须合并，不能合并的差异必须留下；\(a^nb^n\) 与括号配平落在有限状态的射程之外，泵引理把这句话变成可以拿去否定别人的工具。

两节都不碰硬件时序、编译器实现和完整的形式语言层级。第一遍读完，能回答这两个问题就够了：当前的判断需要记住历史吗？若需要，哪些历史差异会改变未来的答案？
